\begin{exercisechunk}
\item \textbf{Loss of plasticity with compatible examples.}
\label[exercise]{ex:l21-compatible-relu}
\exerciseprofile{2}{2}{2}{\exproof\quad\excalculation\quad\exinterpretation}{%
Separate the existence of a common predictor from the ability of gradient updates to reach it.}
Consider $f_{a,b}(x)=\max\{ax+b,0\}$ and the loss
$\tfrac12(f_{a,b}(x)-y)^2$. Initialize $(a,b)=(0,1)$ and apply gradient
descent to each presented example with stepsize $\eta=2/5$.
\begin{enumerate}[(a)]
\item Show that one update on $(x,y)=(1,0)$ gives
$(a,b)=(-2/5,3/5)$. Show that subsequent training only on $(2,1)$ leaves
these parameters unchanged, with loss $1/2$ on every round.
\item Exhibit parameters that fit both $(1,0)$ and $(2,1)$ exactly.
Explain why part (a) is a failure of the update dynamics rather than
incompatibility of the two examples. Also compare the future loss with
that of the prescribed initialization on $(2,1)$.
\item Suppose the same two examples are instead drawn independently
forever, with $\P((X,Y)=(1,0))=q$ and
$\P((X,Y)=(2,1))=1-q$, where $0<q<1$.
Prove that the expected one-round loss converges to $(1-q)/2$,
although a predictor with zero population loss exists.
\end{enumerate}
\begin{hints}
For (c), examine the first occurrence of $(1,0)$ and then track
$a+b$ and $2a+b$. Each example occurs infinitely often almost surely.
\end{hints}
\end{exercisechunk}
